\documentclass[a4paper]{article}

\usepackage[margin=2cm]{geometry}
\usepackage{amssymb}
\usepackage{pdfpages}
\usepackage{graphbox}
\usepackage{tikz}
\begin{document}
\rule{0mm}{14mm}
\begin{center}
  {\Large\bf Values of the standard normal distribution function}
\end{center}

\rule{0mm}{5mm}
\pagestyle{empty}
\thispagestyle{empty}
\begin{table}[ht]
\centering
\begin{tabular}{rllllllllll}
  \hline \rule{0mm}{4mm}
\input{gaussian_table_values.txt}
   \hline
\end{tabular}
\end{table}

\noindent
\(\displaystyle P(Z\leqslant z)=\Phi(z)=\frac{1}{\sqrt{2\pi}}\int_{-\infty}^z\exp\left(-u^2/2\right)\,\mathrm{d}u\)\hfill\hbox{\includegraphics[align=c,width=2.2in]{gaussian_diagram.pdf}}

\noindent For example, to calculate \(P(Z\leqslant 2.93)\), the last
line gives values for \(2.90-2.99\).  To find the value for 2.93,
choose the entry with column heading ``0.03'' to obtain 0.9983.

\begin{tikzpicture}[remember picture,overlay]
  \node at (current page.south east)[anchor = south east]{\includegraphics[width=1in]{AUT-logo-block.pdf}\phantom{\rule[-10mm]{10mm}{10mm}}};
\end{tikzpicture}



\end{document}
